\documentclass[11pt,a4paper]{article}
\usepackage[a4paper,margin=27mm]{geometry}
\usepackage[T1]{fontenc}
\usepackage{lmodern}
\usepackage{amsmath,amssymb,amsthm}
\usepackage{booktabs}
\usepackage{microtype}
\usepackage[hidelinks]{hyperref}
\usepackage{url}
\usepackage{array}
\setlength{\emergencystretch}{3em}

\newtheorem{theorem}{Theorem}[section]
\newtheorem{lemma}[theorem]{Lemma}
\newtheorem{proposition}[theorem]{Proposition}
\newtheorem{corollary}[theorem]{Corollary}
\theoremstyle{definition}
\newtheorem{definition}[theorem]{Definition}
\newtheorem{remark}[theorem]{Remark}
\newtheorem{question}[theorem]{Question}

\newcommand{\F}{\mathbb{F}}
\newcommand{\Fn}{\mathcal{F}}
\newcommand{\Aff}{\mathcal{A}}
\newcommand{\Tg}{\mathcal{T}}
\newcommand{\MC}{\operatorname{MC}}
\newcommand{\rem}{\operatorname{rem}}
\newcommand{\spn}{\operatorname{span}}
\newcommand{\GL}{\operatorname{GL}}
\newcommand{\code}[1]{\texttt{#1}}
\newcommand{\repo}{\url{https://github.com/xennygrimmato/math/tree/3bf15be93674c5382d9fc816922abdc1e70c0303/tools/mc222}}

\title{The multiplicative complexity of $2\times 2$ matrix\\ multiplication over $\F_2$ is seven}
\author{[author]}
\date{26 November 2026}

\begin{document}
\maketitle

\begin{abstract}
Let $c_{ij}=a_{i1}b_{1j}\oplus a_{i2}b_{2j}$ $(i,j\in\{1,2\})$ be the four entries of the product of two $2\times2$ matrices over $\F_2$, viewed as Boolean functions of the eight entries. We show that every circuit over the basis $\{\wedge,\oplus,1\}$ that computes all four of them has at least seven AND gates. The circuit model is unrestricted: an AND gate may multiply any two affine combinations of the inputs and of earlier AND outputs, and intermediate functions may have any degree, with $x^2=x$. Strassen's algorithm gives seven AND gates, so the multiplicative complexity is exactly $7$. The classical lower bounds of Winograd and of Hopcroft and Kerr concern bilinear and arithmetic algorithms and do not apply to this model. The proof is a computer verification. Four short lemmas, proved here, reduce the question to a finite search. Two independent programs then carry out the search: an exact search over spaces of truth tables, and a SAT encoding split into 1080 cubes, each refuted by CaDiCaL 2.1.3 with an LRAT proof that \code{lrat-check} accepts. No part of the argument is formalized in a proof assistant.
\end{abstract}

\section{Introduction}

The multiplicative complexity of a Boolean function, or of a list of Boolean functions, is the least number of AND gates in a circuit over the basis $\{\wedge,\oplus,1\}$ that computes it \cite{BoyarPeraltaPochuev00,BoyarFind18}. XOR gates and the constant $1$ are free. So every input of an AND gate is an affine combination, over $\F_2$, of the circuit inputs and of the outputs of earlier AND gates, and every circuit output is such a combination.

Write the inputs as $x=(a_{11},a_{12},a_{21},a_{22},b_{11},b_{12},b_{21},b_{22})\in\F_2^8$ and let $c_{ij}(x)=a_{i1}b_{1j}+a_{i2}b_{2j}$, with addition in $\F_2$.

\begin{theorem}\label{thm:main}
The multiplicative complexity of the map $\F_2^8\to\F_2^4$, $x\mapsto(c_{11},c_{12},c_{21},c_{22})$, is $7$.
\end{theorem}

The upper bound is Strassen's algorithm \cite{Strassen69}, which has no divisions and integer coefficients and therefore works over $\F_2$ (Proposition~\ref{prop:strassen}). The content of Theorem~\ref{thm:main} is the lower bound: no circuit with six AND gates exists.

\paragraph{Why the classical bound does not apply.} Winograd \cite{Winograd71} proved that seven multiplications are necessary to multiply two $2\times2$ matrices when the entries are indeterminates and the algorithm must produce the product as an identity of polynomials. Hopcroft and Kerr \cite{HopcroftKerr71} proved the bound for algorithms that do not use commutativity, and their argument gives bilinear rank at least $7$ over $\F_2$ \cite{AlderStrassen26}. A Boolean circuit is weaker in what it must achieve and stronger in what it may use. It needs to be correct only on the $256$ points of $\F_2^8$, so it computes functions, not polynomials, and the relation $x^2=x$ holds. An AND gate may take as input an earlier AND output, so intermediate functions can have degree above two, and terms of high degree may later collapse to degree two. The standard passage from general arithmetic circuits to quadratic or bilinear ones, Strassen's homogenization \cite{Strassen73}, does not survive this (Section~\ref{sec:arith}). For one or two quadratic output functions over $\F_2$, Mirwald and Schnorr proved that circuits whose gates have only affine inputs are optimal \cite{MirwaldSchnorr92}, as stated in \cite[\S4]{BoyarFind18}; Boyar and Find state that the corresponding question for more outputs is open \cite[\S4]{BoyarFind18}. The matrix product has four outputs.

\paragraph{What is claimed.} Theorem~\ref{thm:main} rests on Lemmas~\ref{lem:normal}--\ref{lem:helpers}, which are proved below, and on a finite computation, which was carried out twice by programs that share no code (Sections~\ref{sec:exact} and~\ref{sec:sat}). The second computation produces proof objects: for each of $1080$ propositional formulas an LRAT refutation, checked by \code{lrat-check} from the drat-trim repository. That a formula is unsatisfiable implies the theorem only if the formula says what it is meant to say; Section~\ref{sec:trust} lists everything that must be trusted. There is no formalization in Lean or any other proof assistant.

\paragraph{Prior statements.} The bounds $6\le \MC\le 7$ were stated publicly on 4 September 2026 in an unrefereed, machine-generated register, which describes the choice between $6$ and $7$ as open \cite{SDMF149,SDMF180,SDML088}; the lower bound $6$ there comes from a two-dimensional subspace of the outputs. For polynomial multiplication over $\F_2$, Morse proved that the unrestricted multiplicative complexity of multiplying two polynomials with four coefficients is $9$, with a Lean formalization, and that the values for two and three coefficients are $3$ and $6$ \cite{Morse26}; matrices are not treated there. We found no statement of the equality $\MC=7$ for the $2\times2$ matrix product in the literature in a search on 26 November 2026. That search used general web search and arXiv and was not exhaustive; some sources were seen only as search snippets. We do not claim priority beyond this.

\section{The model}\label{sec:model}

Let $\Fn$ be the $\F_2$-algebra of functions $\F_2^8\to\F_2$ with pointwise operations. It is a vector space of dimension $256$, and every $f\in\Fn$ satisfies $f^2=f$. Write $x_1,\dots,x_8$ for the coordinate functions, in the order $a_{11},a_{12},a_{21},a_{22},b_{11},b_{12},b_{21},b_{22}$. Let
\[
\Aff=\spn\{1,x_1,\dots,x_8\},\qquad \Tg=\spn\{c_{11},c_{12},c_{21},c_{22}\}.
\]

\begin{definition}\label{def:circuit}
A \emph{circuit with $k$ gates} is a sequence $g_1,\dots,g_k\in\Fn$ such that $g_i=L_iR_i$ for some $L_i,R_i\in S_{i-1}$, where $S_j=\Aff+\spn\{g_1,\dots,g_j\}$. It \emph{computes} a set of functions $Y$ if $Y\subseteq S_k$. The \emph{multiplicative complexity} $\MC(Y)$ is the least $k$ for which a circuit with $k$ gates computes $Y$.
\end{definition}

A circuit over $\{\wedge,\oplus,1\}$ with $k$ AND gates gives such a sequence: list the AND gates in topological order. Each XOR node computes an element of the current space $S_j$, and NOT is XOR with $1$. Conversely each $L_i$, $R_i$ and each output in $S_k$ is a finite XOR of generators. So Definition~\ref{def:circuit} is the usual multiplicative complexity. Nothing restricts the degree of $g_i$ as a polynomial.

We call a gate $g_i$ a \emph{reducer} if $g_i\in(S_{i-1}+\Tg)\setminus S_{i-1}$ and a \emph{helper} if $g_i\notin S_{i-1}+\Tg$.

\begin{proposition}[Strassen over $\F_2$]\label{prop:strassen}
$\MC(\Tg)\le 7$. Explicitly, with
\begin{align*}
m_1&=(a_{11}+a_{22})(b_{11}+b_{22}), & m_2&=(a_{21}+a_{22})\,b_{11}, & m_3&=a_{11}(b_{12}+b_{22}),\\
m_4&=a_{22}(b_{11}+b_{21}), & m_5&=(a_{11}+a_{12})\,b_{22}, & m_6&=(a_{11}+a_{21})(b_{11}+b_{12}),\\
m_7&=(a_{12}+a_{22})(b_{21}+b_{22}),
\end{align*}
one has $c_{11}=m_1+m_4+m_5+m_7$, $c_{12}=m_3+m_5$, $c_{21}=m_2+m_4$ and $c_{22}=m_1+m_2+m_3+m_6$.
\end{proposition}

\begin{proof}
These are Strassen's identities \cite{Strassen69} with all signs replaced by $+$, which is valid in characteristic $2$. Expanding, for example, $m_3+m_5=a_{11}b_{12}+a_{11}b_{22}+a_{11}b_{22}+a_{12}b_{22}=c_{12}$; the other three are similar. The circuit file \code{results/circuit\_mm222\_7\_p2.txt} contains exactly these seven products, in the order $m_1,m_2,m_4,m_5,m_7,m_3,m_6$, and these four outputs; \code{verify\_circuit.py} evaluates it on all $256$ inputs and accepts it.
\end{proof}

\section{Arithmetic lower bounds and the Boolean model}\label{sec:arith}

Bilinear algorithms are circuits in the sense of Definition~\ref{def:circuit} in which every $L_i$ is a linear form in the $a$'s and every $R_i$ a linear form in the $b$'s. Quadratic algorithms allow every $L_i,R_i\in\Aff$. Hence
\[
\MC(\Tg)\;\le\;\text{(quadratic complexity)}\;\le\;\text{(bilinear rank over }\F_2)=7 .
\]
Lower bounds for general arithmetic circuits over an infinite field are reduced to the quadratic case by homogenization \cite{Strassen73}. One expands each intermediate polynomial $f$ into homogeneous parts $f=f^{(0)}+f^{(1)}+f^{(2)}+\cdots$ and uses that for a product
\[
(fg)^{(2)}=f^{(0)}g^{(2)}+f^{(1)}g^{(1)}+f^{(2)}g^{(0)} ,
\]
so that the degree-two parts of all intermediates are computed by one new product $f^{(1)}g^{(1)}$ per gate. In $\Fn$ there is no grading that is multiplicative in this way. Functions are represented by multilinear polynomials, and the product of two multilinear monomials $x_Sx_{S'}$ is $x_{S\cup S'}$, whose degree $|S\cup S'|$ can be smaller than $|S|+|S'|$. For instance $g=x_1x_2$ has no constant or linear part, yet $g\cdot g=x_1x_2$ and $g\cdot x_1=x_1x_2$. So the quadratic part of a product depends on quadratic and higher parts of the factors, and the induction behind homogenization does not go through. We know of no quadratic map over $\F_2$ whose unrestricted multiplicative complexity is known to be smaller than its quadratic complexity; for one or two outputs there is none \cite{MirwaldSchnorr92}. The point is only that no known general argument transfers the bound $7$ to the Boolean model.

\section{Reductions}\label{sec:lemmas}

This section proves the four lemmas used by the computations, and one more that explains the running times. The table at the end records which method uses which lemma.

For a circuit, write each $L_i$, $R_i$ as an $\F_2$-combination of the \emph{generators} $1,x_1,\dots,x_8$, $g_1,\dots,g_{i-1}$. When the generators are linearly independent, the coefficient vector of an element of $S_{i-1}$ is unique.

\begin{lemma}[normal form]\label{lem:normal}
\begin{enumerate}
\item[(a)] If in a circuit one replaces $g_i$ by $g_i'=g_i+s$ with $s\in S_{i-1}$, and $g_i'=L_i'R_i'$ with $L_i',R_i'\in S_{i-1}$, then the result is a circuit with the same spaces $S_j$ for all $j$.
\item[(b)] For $u,v\in S_{i-1}$: $(u+1)v-uv\in S_{i-1}$, and $u'v'-uv\in S_{i-1}$ for every basis $(u',v')$ of $\spn\{u,v\}$.
\item[(c)] If $\MC(\Tg)\le k$, then some circuit with exactly $k$ gates computes $\Tg$ in which, for each $i$, $L_i$ and $R_i$ have constant coefficient $0$ and their coefficient vectors are the reduced row-echelon basis of a two-dimensional space (pivot of $L_i$ before pivot of $R_i$, each vector zero at the other's pivot).
\item[(d)] In a circuit with the least number of gates computing $\Tg$, every gate satisfies $g_i\notin S_{i-1}$, so $\dim S_i=9+i$.
\end{enumerate}
\end{lemma}

\begin{proof}
(a) $S_{i-1}$ is unchanged and $S_{i-1}+\spn\{g_i'\}=S_{i-1}+\spn\{g_i\}$. By induction on $j\ge i$ the spaces $S_j$ agree, so every later factor $L_j,R_j$ still lies in $S_{j-1}$ and every later gate is still a product of two elements of the previous space.
(b) $(u+1)v=uv+v$. The bases of a two-dimensional space over $\F_2$ are the ordered pairs of distinct nonzero elements of $\{u,v,u+v\}$, and $u(u+v)=uv+u$ and $(u+v)v=uv+v$ because $u^2=u$ and $v^2=v$; the products for the other bases follow by symmetry.
(d) If $g_i\in S_{i-1}$, delete it. The spaces $S_j$ for $j>i$ do not change, and later factors, expressed through the remaining generators, still lie in the previous spaces. This gives a shorter circuit.
(c) Take a circuit with the least number $m\le k$ of gates. By (d) the generators are linearly independent. If the coefficient vectors of $L_i$ and $R_i$ (constant dropped) were zero or equal, then $g_i$ would be a product of two constants or of $u+\epsilon$ and $u+\epsilon'$ with $\epsilon,\epsilon'\in\{0,1\}$, hence $g_i\in S_{i-1}$, contradicting (d). So they span a two-dimensional space, and by (b) and (a) we may replace $(L_i,R_i)$ by its reduced echelon basis without constants. Finally append $k-m$ gates $x_1\cdot x_2$, whose coefficient vectors $e_1,e_2$ are in reduced echelon form. The appended gates do not shrink $S_k$.
\end{proof}

Define $\rem(S)=\dim(S+\Tg)-\dim S$ for subspaces $S\supseteq\Aff$ of $\Fn$.

\begin{lemma}[counting]\label{lem:count}
$\rem(\Aff)=4$. Let $S\supseteq\Aff$ and $g\notin S$. If $g\in S+\Tg$ then $\rem(S+\spn\{g\})=\rem(S)-1$ and $S+\spn\{g\}=S+\spn\{t\}$ for some $t\in\Tg\setminus S$. Otherwise $\rem(S+\spn\{g\})=\rem(S)$. Consequently, in a circuit with the least number $m$ of gates computing $\Tg$, exactly $4$ gates are reducers and $m-4$ are helpers, and a state $S_j$ with $\rem(S_j)>k-j$ cannot be extended to a $k$-gate circuit computing $\Tg$.
\end{lemma}

\begin{proof}
The multilinear monomials form a basis of $\Fn$. The four functions $c_{ij}$ are sums of the eight distinct monomials $a_{ik}b_{kj}$, each used once, so they are linearly independent modulo $\Aff=\spn\{1,x_1,\dots,x_8\}$, and $\rem(\Aff)=4$. If $g=s+t$ with $s\in S$, $t\in\Tg$, then $S+\spn\{g\}=S+\spn\{t\}$, $S+\spn\{g\}+\Tg=S+\Tg$, and the dimension of $S$ grows by one. If $g\notin S+\Tg$, both $\dim S$ and $\dim(S+\Tg)$ grow by one. A circuit computes $\Tg$ exactly when $\rem(S_k)=0$. By Lemma~\ref{lem:normal}(d) every gate of a shortest circuit is a reducer or a helper, and each reducer lowers $\rem$ by exactly one.
\end{proof}

Let $G$ be the group of linear maps $\F_2^8\to\F_2^8$ generated by $(A,B)\mapsto(PAQ^{-1},QBR^{-1})$ for $P,Q,R\in\GL_2(\F_2)$ and by $(A,B)\mapsto(B^{\mathsf T},A^{\mathsf T})$. The map $(P,Q,R)\mapsto$ (first kind) is injective, since $PAQ^{-1}=A$ for all $A$ forces $P=Q=I$, and the transposition is not of the first kind, so $|G|=2\cdot 6^3=432$. A plane is a two-dimensional subspace of $\F_2^8$; there are $255\cdot254/6=10795$ of them. Identify $a\in\F_2^8$ with the linear form $\ell_a(x)=\sum_k a_kx_k$; then $\ell_a\circ\sigma=\ell_{\sigma^{\mathsf T}a}$, and $G$ acts on planes.

\begin{lemma}[symmetry]\label{lem:sym}
For $\sigma\in G$, the map $f\mapsto f\circ\sigma$ is an algebra automorphism of $\Fn$ with $\Aff\circ\sigma=\Aff$ and $\Tg\circ\sigma=\Tg$. If $g_1,\dots,g_m$ is a circuit computing $\Tg$, so is $g_1\circ\sigma,\dots,g_m\circ\sigma$, with spaces $S_j\circ\sigma$. In a shortest circuit $g_1=L_1R_1$ with $L_1,R_1\in\Aff$ whose linear parts $\ell_a,\ell_b$ are independent; replacing the circuit by its image under a suitable $\sigma$ and applying Lemma~\ref{lem:normal}, one may assume that $\spn\{a,b\}$ is any prescribed plane in the $G$-orbit of the original one, and that $(L_1,R_1)$ is its reduced echelon basis.
\end{lemma}

\begin{proof}
Composition with a bijection preserves sums and products. Since $\sigma$ is linear and invertible, $\Aff\circ\sigma=\Aff$. Under the first kind of map $AB\mapsto P(AB)R^{-1}$, whose entries are linear combinations of the entries $c_{ij}$; under the second, $B^{\mathsf T}A^{\mathsf T}=(AB)^{\mathsf T}$, whose entries are a permutation of the $c_{ij}$. So $\Tg\circ\sigma\subseteq\Tg$, with equality by invertibility. Then $(g_i\circ\sigma)=(L_i\circ\sigma)(R_i\circ\sigma)$ with factors in $S_{i-1}\circ\sigma$, and $\Tg=\Tg\circ\sigma\subseteq S_m\circ\sigma$. The first factors become $\ell_{\sigma^{\mathsf T}a}$ and $\ell_{\sigma^{\mathsf T}b}$ plus constants. The independence of $a,b$ is Lemma~\ref{lem:normal}(c)--(d); constants and the choice of basis are removed by Lemma~\ref{lem:normal}(a)--(b), which changes no space.
\end{proof}

\begin{proposition}[computed]\label{prop:orbits}
The $10795$ planes fall into $72$ orbits under $G$, of sizes $4,6^2,9,12^2,18^3,27^2,36^5,54^5,72^2,108^{19},216^{23},432^7$ (exponents are multiplicities). This was computed twice: by \code{orbits.py}, which builds all $432$ elements of $G$, checks that each is linear, invertible and preserves $\Aff+\Tg$, and writes \code{plane\_orbits.txt}; and by \code{src/main.rs}, which closes each plane under six generators of $G$ and the transposition after checking the same properties on truth tables. Both choose the lexicographically least pair $(a,b)$ with $a<b<a\oplus b$ (as $8$-bit masks, bit $k-1$ for $x_k$) as representative, and both give the same $72$ representatives and the same sizes.
\end{proposition}

\begin{lemma}[the first gate is a helper]\label{lem:first}
If $a,b\in\F_2^8$ are linearly independent, then $\ell_a\ell_b\notin\Aff+\Tg$.
\end{lemma}

\begin{proof}
For $f\in\Fn$ put $B_f(y,z)=f(y+z)+f(y)+f(z)+f(0)$. The map $f\mapsto B_f$ is linear, vanishes on $\Aff$, and sends $\ell_a\ell_b$ to the symmetric form with matrix $ab^{\mathsf T}+ba^{\mathsf T}$, of rank $2$. For $t=\sum\mu_{ij}c_{ij}\ne0$ we have $t=\sum_{i,j,k}\mu_{ij}a_{ik}b_{kj}$, so $B_t$ has the matrix $\begin{pmatrix}0&M\\M^{\mathsf T}&0\end{pmatrix}$ in the $(A,B)$ coordinates, where $M_{(i,k),(k',j)}=\mu_{ij}\delta_{kk'}$. After reordering, $M$ is block diagonal with two blocks equal to $\mu$, so $\operatorname{rank}B_t=2\operatorname{rank}M=4\operatorname{rank}\mu\in\{4,8\}$. Hence $B_{\ell_a\ell_b}\neq B_t$ for every $t\in\Tg$.
\end{proof}

\begin{lemma}[helper sets]\label{lem:helpers}
Suppose $\MC(\Tg)\le 6$. Then for one of the $72$ orbit representatives $(\alpha,\beta)$ there is a circuit $g_1,\dots,g_6$ in the normal form of Lemma~\ref{lem:normal}(c), with $(L_1,R_1)$ the reduced echelon basis of $\spn\{\alpha,\beta\}$, and a set $H\subseteq\{1,\dots,6\}$ with $|H|=2$, such that every $i\notin H$ satisfies
\[
g_i\in W_i:=\Aff+\Tg+\spn\{g_j: j\in H,\ j<i\}.
\]
Moreover $1\in H$.
\end{lemma}

\begin{proof}
Take a shortest circuit, with $m\le 6$ gates; $m\ge4$ by Lemma~\ref{lem:count}. By Lemma~\ref{lem:sym} we may assume that its first plane is a representative and, by Lemma~\ref{lem:normal}(a)--(c), that it is in normal form. These operations keep the number of gates, so the circuit is still shortest, and they change no space $S_j$, so each gate stays a reducer or a helper. By Lemma~\ref{lem:count} there are $m-4$ helpers. Append $6-m$ padding gates $x_1x_2$ and let $H$ consist of the helpers and the padding positions; then $|H|=(m-4)+(6-m)=2$. We show $S_{i-1}\subseteq W_i$ for $i\le m$ by induction on $i$. For $i=1$, $S_0=\Aff$. For $j<i$: if $j\in H$ then $g_j\in W_i$ by definition; if $j\notin H$ then $g_j$ is a reducer, so $g_j\in S_{j-1}+\Tg\subseteq W_j+\Tg=W_j\subseteq W_i$. Hence for each reducer $g_i\in S_{i-1}+\Tg\subseteq W_i$. Padding positions lie in $H$. Finally $g_1$ is a helper by Lemma~\ref{lem:first}, so $1\in H$.
\end{proof}

\begin{table}[h]
\centering
\small
\begin{tabular}{>{\raggedright}p{0.22\textwidth}>{\raggedright}p{0.36\textwidth}>{\raggedright\arraybackslash}p{0.32\textwidth}}
\toprule
statement & exact search (A) & SAT with proofs (B)\\
\midrule
Lemma~\ref{lem:normal} & no constants; one product per plane; $\dim S_j=9+j$ & no constants; echelon pair; padding\\
Lemma~\ref{lem:count} & pruning $\rem>b$; reducer successors $S+\spn\{t\}$ & only through Lemma~\ref{lem:helpers}\\
Lemma~\ref{lem:sym}, Prop.~\ref{prop:orbits} & $72$ roots (Rust orbits) & $72$ roots (Python orbits)\\
Lemma~\ref{lem:first} & not used & not used (explains timings)\\
Lemma~\ref{lem:helpers} & not used & the $15$ cube constraints per root\\
\bottomrule
\end{tabular}
\caption{Use of the reductions. Both methods depend on Lemmas~\ref{lem:normal}, \ref{lem:count} and~\ref{lem:sym} and on the orbit computation, which each method performs with its own code.}
\label{tab:use}
\end{table}

\section{Method A: exact search over spaces}\label{sec:exact}

The program \code{src/main.rs} (813 lines of Rust, no dependencies) decides whether a $k$-gate circuit computes $\Tg$ by depth-first search over the spaces $S_j$, stored as $256$-bit truth tables. By Lemma~\ref{lem:normal}(d) it suffices to consider gates with $g_j\notin S_{j-1}$, so $\dim S_j=9+j$ and the remaining budget is $b=k-j$, a function of $S_j$. A state is identified by the fully reduced echelon basis of $S_j$ modulo $\Aff$, which determines $S_j$. Failed states are stored with this full key in a hash set, so two different states are never identified. The cache depends only on the state and is therefore sound; the current code clears it at the start of each root.

At a state $S$ the program computes $\rem(S)$ and stops if $\rem(S)>b$ (Lemma~\ref{lem:count}). Its successors are the spaces $S+\spn\{uv\}$ with $u,v\in S$ and $uv\notin S$; by Lemma~\ref{lem:normal}(b) constants may be dropped from $u$ and $v$. These are produced in two classes.

\emph{Reducers.} For fixed $u$, the map $v\mapsto uv$ is linear. Gaussian elimination of the products of $u$ with a basis of $S$, reduced modulo $S+\Tg$ and tagged with their coordinates in $(S+\Tg)/S$, yields the subspace of classes $t\in(S+\Tg)/S$ that some $uv$ hits. The union over all $2^{\dim S-1}-1$ nonzero $u$, enumerated in Gray-code order, is the set of reachable classes; each nonzero class $t$ gives the successor $S+\spn\{t\}$ (Lemma~\ref{lem:count}).

\emph{Helpers}, only when $\rem(S)<b$. The program enumerates one pair $u<v<u\oplus v$ per plane of $S/\spn\{1\}$, which suffices by Lemma~\ref{lem:normal}(b), keeps the products with $uv\notin S+\Tg$, and removes duplicates by the reduced residue of $uv$ modulo $S$. Since $S+\spn\{w\}\setminus S=w+S$, this residue determines the successor, and the deduplication is exact.

The roots are the $72$ states $\Aff+\spn\{\ell_\alpha\ell_\beta\}$ (Lemma~\ref{lem:sym}, Proposition~\ref{prop:orbits}).

\paragraph{Result for $k=6$.} Command: \code{target/release/mc222 mm222 6 -{}-threads 4}. All $72$ roots report no circuit. In the re-run of 26 November 2026 the search visited $3{,}547{,}604$ states in $462.1$\,s wall time on four threads. Sixty-nine roots visit between $42{,}196$ and $42{,}358$ states; roots $4$, $5$ and $36$ visit $211{,}365$, $211{,}267$ and $211{,}169$. Per-root counts are deterministic: we re-ran roots $68$--$71$ on one thread and obtained $42{,}196$ states each, as in the four-thread run. The log committed in \code{results/exact\_k6\_orbits.log} records the same verdicts with $3{,}546{,}906$ states and $396.8$\,s. Its per-root counts agree with the re-run for the first root of each thread and are slightly smaller for later roots. This is what a build that kept the cache across the roots of one thread would produce, which is sound; we have not identified that build. For the same reason the committed description gives $10{,}943$ states for $k=5$ without orbit reduction, where the current code gives $10{,}975$.

\section{Method B: SAT with checked refutations}\label{sec:sat}

The program \code{gen\_cnf.py} (246 lines of Python, written separately from the Rust code) writes a CNF formula over truth tables at all $256$ input points. Encodings of this kind, with one copy of the circuit per input point, have been used to compute the multiplicative complexity of S-boxes \cite{Stoffelen16}. Gate positions are numbered from $1$ here and from $0$ in the code.

\emph{Gates.} Gate $i$ has coefficient variables $\lambda_i,\rho_i\in\F_2^{7+i}$ over $x_1,\dots,x_8,g_1,\dots,g_{i-1}$, with constraints stating that both are nonzero, that the first nonzero position of $\lambda_i$ precedes that of $\rho_i$, and that each is zero at the other's first nonzero position. This is the normal form of Lemma~\ref{lem:normal}(c). At each point $p$ the value of the linear part is built by the recursion $X(p)=X(p-2^h)\oplus s_h$, $h$ the highest set bit of $p$, the contributions of earlier gates are added, and $g_i(p)=L_i(p)\wedge R_i(p)$, all by Tseitin clauses with constant folding.

\emph{Outputs.} Each $c_{ij}$ gets free coefficients for $1$, $x_1,\dots,x_8$ and $g_1,\dots,g_6$, and the combination must equal $c_{ij}(p)$ at all $256$ points.

\emph{Cubes.} For a representative $(\alpha,\beta)$, gate $1$ is fixed to the reduced echelon basis of $\spn\{\alpha,\beta\}$. For a set $H$ of two positions, each gate $i\notin H$ gets fresh variables $e,\ell_1,\dots,\ell_8,\mu_1,\dots,\mu_4$ and $\gamma_j$ $(j\in H,\ j<i)$ with the constraint $g_i=e+\sum\ell_kx_k+\sum\mu_tc_t+\sum\gamma_jg_j$ at all $256$ points, that is, $g_i\in W_i$ (Lemma~\ref{lem:helpers}).

There are $72\cdot\binom{6}{2}=1080$ cubes. The cubes with $1\notin H$ are unsatisfiable by Lemma~\ref{lem:first}; they were run anyway. A typical formula, root $0$ with $H=\{1,4\}$, has $42{,}036$ variables and $152{,}710$ clauses ($2.89$\,MB).

\paragraph{Result for $k=6$.} Command: \code{python3 run\_sat.py <cadical> <lrat-check> 3 6}, then \code{python3 summarize\_sat.py out/sat\_k6/results.tsv}. For each cube, \code{run\_sat.py} generates the formula, runs CaDiCaL 2.1.3 \cite{CaDiCaL2} (\code{-q -{}-lrat=true -{}-binary=false}), checks the LRAT proof \cite{CruzFilipe17} with \code{lrat-check} from drat-trim \cite{Wetzler14} (commit \code{2e3b2dc0ecf938add}\allowbreak\code{bd779d42877b6ed69d9a985}), records sizes and SHA-256 hashes of formula and proof, and deletes both. All $1080$ cubes are unsatisfiable and all $1080$ proofs are reported \code{c VERIFIED}. The formulas total $3.05$\,GB ($2.68$--$2.96$\,MB each) and the proofs $5.41$\,GB ($0.40$--$54.05$\,MB each; the largest is root $65$, $H=\{1,2\}$). In the committed record solving took $0.33$\,h CPU in total (at most $17.2$\,s per cube) and checking $0.08$\,h; the re-run took $0.59$\,h and $0.13$\,h on a more loaded machine. Most of the work is in $H=\{1,2\}$, $\{1,3\}$ and $\{1,4\}$ ($685$\,s, $263$\,s and $116$\,s in the committed record); the $720$ cubes with $1\notin H$ took $56$\,s together.

\paragraph{Proof objects.} An LRAT check certifies that the formula file it was given has no satisfying assignment, provided the checker is correct; it does not certify that the formula expresses the intended statement. The proofs are not stored in the repository. The file \code{results/sat\_k6\_results.tsv} lists, for each cube, the representative, $H$, verdict, check result, times, and the size and SHA-256 hash of the formula and of the proof. CaDiCaL is deterministic, so one script regenerates all of them: \code{run\_all.sh} builds the pinned solver and checker and reruns everything. In the independent re-run every one of the $1080$ formula hashes and $1080$ proof hashes was identical to the committed record. We also regenerated root $0$ with $H=\{1,4\}$ separately and obtained the recorded hashes \code{8561c7e2\ldots} (formula) and \code{3854bd00\ldots} (proof), and a \code{VERIFIED} check.

\begin{proof}[Proof of Theorem~\ref{thm:main}]
The upper bound is Proposition~\ref{prop:strassen}. Suppose $\MC(\Tg)\le6$. By Lemma~\ref{lem:helpers} some cube of Method B has a circuit satisfying all its constraints, and the formula of that cube is then satisfiable. All $1080$ formulas are unsatisfiable. Independently, by Lemmas~\ref{lem:normal}, \ref{lem:count} and~\ref{lem:sym} some shortest circuit, with $m\le6$ gates, has its first space $S_1$ equal to one of the $72$ roots of Method A. Its spaces $S_1\subset\dots\subset S_m$ form a path of successors in the sense of Section~\ref{sec:exact}, along which the pruning bound of Lemma~\ref{lem:count} never applies, and the search stops with success at any state with $\rem=0$, so no padding is needed. Method A finds no such path.
\end{proof}

\section{Validation and reproduction}\label{sec:validation}

Both programs were run on small problems with known answers. Each circuit found by either program, and each model of a satisfiable formula decoded into a circuit, was checked by \code{verify\_circuit.py}, which re-defines the targets and evaluates the circuit at every input. Table~\ref{tab:val} gives the results of our re-run of the validation part of \code{run\_all.sh}; the SAT runs here do not write proofs.

\begin{table}[h]
\centering
\footnotesize
\begin{tabular}{>{\raggedright}p{0.2\textwidth}>{\raggedright}p{0.07\textwidth}>{\raggedright}p{0.3\textwidth}>{\raggedright\arraybackslash}p{0.3\textwidth}}
\toprule
problem & value & (A) exact search & (B) SAT\\
\midrule
$\F_4$ multiplication, $4$ inputs & $3$ & $k=2$ none, $k=3$ found & $k=2$ UNSAT, $k=3$ SAT\\
$2$-term polynomials, $4$ inputs & $3$ \cite{Morse26} & $k=2$ none, $k=3$ found & $k=2$ UNSAT, $k=3$ SAT\\
$3$-term polynomials, $6$ inputs & $6$ \cite{Morse26} & $k=5$ none, $k=6$ found & $k=5$ UNSAT ($49$\,s), $k=6$ SAT ($8$\,s)\\
$\langle2,2,1\rangle$, $6$ inputs & $4$ & $k=3$ none, $k=4$ found & $k=3$ UNSAT, $k=4$ SAT\\
$\langle2,2,2\rangle$, $k=5$ & & none ($72$ roots, $75$ states; all $10795$ planes, $10975$ states) & root $0$ UNSAT ($5$\,s)\\
$\langle2,2,2\rangle$, $k=7$ & & found with $m_1,m_2$ fixed ($0.7$\,s) & $H=\{1,2,3\}$, no root: SAT ($73$\,s)\\
\bottomrule
\end{tabular}
\caption{Validation runs. Times are wall-clock on one core for (B).}
\label{tab:val}
\end{table}

For $\F_4$ multiplication the lower bound $3$ follows from Lemma~\ref{lem:count} and the argument of Lemma~\ref{lem:first}: with two gates and $\rem=2$ both gates would be reducers, and the first, a product of two linear forms, would lie in $\Aff+\Tg$, while the three nonzero target combinations have alternating rank $4$. For $\langle2,2,1\rangle$ the value $4$ is the one found by our programs and by an earlier SAT computation in our notebook; we did not locate a published source for its unrestricted multiplicative complexity. The SAT run for $k=7$ uses $H=\{1,2,3\}$, the analogue of Lemma~\ref{lem:helpers} with three helpers, and fixes no gate; the decoded circuit is correct on all inputs, which shows that the encoding admits genuine circuits.

\paragraph{Reproduction.} \code{./run\_all.sh [ext\_dir]} clones and builds CaDiCaL at tag \code{rel-2.1.3} and drat-trim at the commit above, builds the Rust program, runs the validation cases, recomputes the orbits, runs Method A for $k=6$ and Method B for all $1080$ cubes, and summarizes. Development versions of CaDiCaL should not be used: version 3.1.0 printed a partial model that the decoder misread.

\section{What is trusted}\label{sec:trust}

\emph{Both methods.} The model (Definition~\ref{def:circuit}) and Lemmas~\ref{lem:normal}--\ref{lem:sym}, proved above. The orbit computation, done twice by independent code with identical results (Proposition~\ref{prop:orbits}). The target definitions, which are written independently in each of \code{main.rs}, \code{gen\_cnf.py}, \code{orbits.py} and \code{verify\_circuit.py}.

\emph{Method A.} The correctness of \code{src/main.rs}: exhaustive generation of successors, the exact state key and the pruning. The Rust compiler (rustc 1.97.0 on the machine of the re-run). Consistency checks: according to the repository's README, running a non-representative member of an orbit as root gave the same state count for six orbits (we did not repeat this); and the search over all $10795$ first gates agrees in verdict with the orbit-reduced search at $k=5$.

\emph{Method B.} The faithfulness of \code{gen\_cnf.py}: if a circuit as in Lemma~\ref{lem:helpers} exists for a cube, the cube's formula is satisfiable. We read the encoding against Section~\ref{sec:sat}. Evidence beyond reading is the validation table: the known satisfiable and unsatisfiable cases come out right, and decoded models are correct circuits. Lemma~\ref{lem:helpers}. The Python orbit file. The checker \code{lrat-check} at the pinned commit, a C program that is not itself verified. CaDiCaL is not trusted.

\emph{Not trusted anywhere.} An earlier search, \code{mc\_span.c}, which first reported that all $72$ orbits fail at $k=6$. It identified states by a $128$-bit hash and ignored collisions, and it produced no proof object. Methods A and B were written without reading it.

\emph{Formalization.} No proof-assistant formalization exists. One would need: the model and Lemmas~\ref{lem:normal}--\ref{lem:helpers}; the orbit computation; a proof that each generated formula is satisfiable whenever its cube contains a circuit; and a verified LRAT checker run on $5.41$\,GB of proofs, such as the certified checkers of \cite{CruzFilipe17}. Alternatively one could certify Method A's search tree.

\section{Related small cases}\label{sec:related}

\begin{remark}[$\langle2,2,1\rangle$]
The matrix-vector product $\F_2^{2\times2}\times\F_2^2\to\F_2^2$ has multiplicative complexity $4$, equal to its bilinear rank \cite{HopcroftKerr71}. Status: computed by both programs (Table~\ref{tab:val}), without proof objects.
\end{remark}

\begin{remark}[$\langle1,2,3\rangle$]
For the product of a $1\times2$ and a $2\times3$ matrix, the multiplicative complexity is $5$ or $6$. The upper bound $6$ is the naive algorithm. The lower bound $5$ rests only on the earlier program \code{mc\_span.c} (all $10795$ first planes refuted at $k=4$), which is not trusted here. The case $k=5$ is open.
\end{remark}

\begin{remark}[quadratic model]
If every gate must multiply two affine functions of the inputs, with $x^2=x$ and with $a$- and $b$-entries allowed in one factor, at least $7$ gates are needed. This is a special case of Theorem~\ref{thm:main}. It was also obtained earlier by a separate enumeration of decomposable alternating forms, which we have not re-checked.
\end{remark}

\section{Open questions}\label{sec:open}

\begin{question}
Is the multiplicative complexity of $\langle1,2,3\rangle$ over $\F_2$ equal to $5$ or $6$?
\end{question}

\begin{question}
What is the least number of multiplications for the $2\times2$ product over $\F_3$ or $\F_4$ in the corresponding model, where gates multiply affine combinations of earlier values and $x^q=x$?
\end{question}

\begin{question}
What is the multiplicative complexity of $\langle2,2,3\rangle$ over $\F_2$? Hopcroft and Kerr give an algorithm with $11$ multiplications and prove that $11$ are needed without commutativity \cite{HopcroftKerr71}.
\end{question}

\begin{question}
Does the multiplicative complexity equal the bilinear rank over $\F_2$ for all small matrix formats? No quadratic map over $\F_2$ is known to us whose multiplicative complexity is smaller than its quadratic complexity.
\end{question}

\begin{question}
Is there a human-readable proof of Theorem~\ref{thm:main}, for example along the lines of Morse's structural argument for polynomial multiplication \cite{Morse26}?
\end{question}

\section{How this paper was produced}\label{sec:provenance}

The searches, the proofs of the lemmas and the text of this paper were produced by AI systems. A first search found no six-gate circuit. A second AI-run pass then re-implemented the search independently, as Methods A and B, without reading the first program, and the whole pipeline was re-run once more with identical verdicts and hashes. No human mathematician has checked the proofs or the code as of the date of writing. The code, the per-cube records and a seven-gate circuit are in the folder \code{tools/mc222} of the repository at commit \code{3bf15be93674c5382d9fc816922abdc1e70c0303}:
\begin{center}\small \repo\end{center}

\bibliographystyle{plain}
\bibliography{refs}

\end{document}
